%%%%%%%%%%%%%%%%%%%%%%%%%%%%%%%%%%%%%%%%%%%%%%%%%%%%%%%%%%%%
% 13.6 CONTROL THEORY
% The Composition of Advanced Mathematics
%%%%%%%%%%%%%%%%%%%%%%%%%%%%%%%%%%%%%%%%%%%%%%%%%%%%%%%%%%%%

\section{Control Theory}
\label{sec:control-theory}

\prerequisite{
Linear algebra, differential equations, dynamical systems, Laplace
transforms, optimization, numerical methods, probability, mathematical
modeling, and mathematical physics.
}

\learningobjectives{
By the end of this section, the reader should be able to:
\begin{itemize}
    \item explain the purpose of control theory;
    \item distinguish open-loop and closed-loop control;
    \item identify plant, controller, actuator, sensor, input, output,
          disturbance, and reference signals;
    \item formulate state-space models;
    \item distinguish state variables from outputs;
    \item understand linear time-invariant systems;
    \item solve linear state equations using matrix exponentials;
    \item understand equilibrium points;
    \item linearize nonlinear systems;
    \item understand transfer functions;
    \item connect state-space and transfer-function representations;
    \item understand poles and zeros;
    \item recognize stability from poles;
    \item understand asymptotic, Lyapunov, and exponential stability;
    \item analyze stability using eigenvalues;
    \item recognize marginal stability;
    \item understand BIBO stability;
    \item construct feedback laws;
    \item understand proportional control;
    \item understand integral control;
    \item understand derivative control;
    \item formulate PID control;
    \item understand steady-state error;
    \item recognize overshoot, rise time, settling time, and damping;
    \item understand characteristic equations;
    \item understand root-locus ideas conceptually;
    \item recognize frequency response;
    \item understand Bode plots conceptually;
    \item recognize gain and phase margins;
    \item understand controllability;
    \item construct controllability matrices;
    \item understand observability;
    \item construct observability matrices;
    \item understand state feedback;
    \item recognize pole placement;
    \item understand observers;
    \item formulate Luenberger observers;
    \item understand separation between estimation and control conceptually;
    \item recognize canonical forms;
    \item understand Lyapunov functions;
    \item verify stability using Lyapunov methods;
    \item recognize Lyapunov equations for linear systems;
    \item understand nonlinear feedback conceptually;
    \item recognize feedback linearization conceptually;
    \item understand optimal control;
    \item formulate quadratic cost functions;
    \item understand linear-quadratic regulation;
    \item recognize Riccati equations;
    \item understand finite- and infinite-horizon control;
    \item recognize Pontryagin's Maximum Principle conceptually;
    \item understand Hamilton--Jacobi--Bellman equations conceptually;
    \item recognize model predictive control;
    \item understand constraints in control;
    \item recognize robust control;
    \item understand uncertainty and model mismatch;
    \item recognize sensitivity functions conceptually;
    \item understand disturbance rejection;
    \item understand tracking;
    \item recognize feedforward control;
    \item understand stochastic control conceptually;
    \item recognize Kalman filtering;
    \item understand state estimation under noise;
    \item recognize discrete-time control systems;
    \item analyze discrete-time stability;
    \item understand digital control;
    \item recognize sampling effects;
    \item understand delay in control systems;
    \item recognize saturation and actuator constraints;
    \item understand anti-windup conceptually;
    \item recognize nonlinear control challenges;
    \item understand adaptive control conceptually;
    \item recognize system identification;
    \item understand controllability and observability Gramians conceptually;
    \item recognize balancing and model reduction conceptually;
    \item understand networked and distributed control conceptually;
    \item connect control theory with optimization, dynamical systems,
          probability, robotics, aerospace, process systems, and finance.
\end{itemize}
}

%%%%%%%%%%%%%%%%%%%%%%%%%%%%%%%%%%%%%%%%%%%%%%%%%%%%%%%%%%%%
\subsection{What Is Control Theory?}
\label{subsec:what-is-control-theory}
%%%%%%%%%%%%%%%%%%%%%%%%%%%%%%%%%%%%%%%%%%%%%%%%%%%%%%%%%%%%

Control theory studies how inputs can be chosen to influence the behavior of
dynamical systems.

The general objective is to make a system:

\[
\boxed{
\text{stable},
}
\]

\[
\boxed{
\text{accurate},
}
\]

\[
\boxed{
\text{fast},
}
\]

\[
\boxed{
\text{robust},
}
\]

or

\[
\boxed{
\text{optimal}
}
\]

according to a desired performance criterion.

A controlled system may be mechanical, electrical, biological, economic, or
computational.

%%%%%%%%%%%%%%%%%%%%%%%%%%%%%%%%%%%%%%%%%%%%%%%%%%%%%%%%%%%%
\subsection{Basic Control-System Components}
\label{subsec:control-system-components}
%%%%%%%%%%%%%%%%%%%%%%%%%%%%%%%%%%%%%%%%%%%%%%%%%%%%%%%%%%%%

A control system may contain:

\[
\boxed{
\text{plant}
}
\]

for the physical or mathematical system being controlled;

\[
\boxed{
\text{controller}
}
\]

for the rule selecting inputs;

\[
\boxed{
\text{actuator}
}
\]

for applying control actions;

\[
\boxed{
\text{sensor}
}
\]

for measuring the system;

and

\[
\boxed{
\text{reference}
}
\]

for the desired behavior.

%%%%%%%%%%%%%%%%%%%%%%%%%%%%%%%%%%%%%%%%%%%%%%%%%%%%%%%%%%%%
\subsection{Open-Loop Control}
\label{subsec:open-loop-control}
%%%%%%%%%%%%%%%%%%%%%%%%%%%%%%%%%%%%%%%%%%%%%%%%%%%%%%%%%%%%

In open-loop control, the control input does not depend on the measured
current output.

Schematically,

\[
\boxed{
r(t)
\longrightarrow
u(t)
\longrightarrow
\text{plant}
\longrightarrow
y(t).
}
\]

The controller does not automatically correct deviations caused by
disturbances or model error.

%%%%%%%%%%%%%%%%%%%%%%%%%%%%%%%%%%%%%%%%%%%%%%%%%%%%%%%%%%%%
\subsection{Closed-Loop Control}
\label{subsec:closed-loop-control}
%%%%%%%%%%%%%%%%%%%%%%%%%%%%%%%%%%%%%%%%%%%%%%%%%%%%%%%%%%%%

Closed-loop control uses feedback.

A measured output is compared with a reference:

\[
\boxed{
e(t)
=
r(t)-y(t).
}
\]

The controller then chooses

\[
\boxed{
u(t)
=
\mathcal{K}[e].
}
\]

Feedback can reduce error and reject disturbances.

%%%%%%%%%%%%%%%%%%%%%%%%%%%%%%%%%%%%%%%%%%%%%%%%%%%%%%%%%%%%
\subsection{Negative Feedback}
\label{subsec:negative-feedback-control}
%%%%%%%%%%%%%%%%%%%%%%%%%%%%%%%%%%%%%%%%%%%%%%%%%%%%%%%%%%%%

Negative feedback acts to oppose deviations from the desired state.

For a simple proportional law,

\[
\boxed{
u
=
K(r-y),
}
\]

where \(K>0\).

If \(y\) becomes too large relative to \(r\), the control action changes in
the direction that reduces \(y\).

%%%%%%%%%%%%%%%%%%%%%%%%%%%%%%%%%%%%%%%%%%%%%%%%%%%%%%%%%%%%
\subsection{State-Space Models}
\label{subsec:control-state-space}
%%%%%%%%%%%%%%%%%%%%%%%%%%%%%%%%%%%%%%%%%%%%%%%%%%%%%%%%%%%%

A general nonlinear state-space model is

\[
\boxed{
\dot{\mathbf{x}}
=
f(\mathbf{x},\mathbf{u},t),
}
\]

\[
\boxed{
\mathbf{y}
=
h(\mathbf{x},\mathbf{u},t).
}
\]

Here:

\[
\mathbf{x}
=
\text{state vector},
\]

\[
\mathbf{u}
=
\text{control input},
\]

and

\[
\mathbf{y}
=
\text{measured output}.
\]

%%%%%%%%%%%%%%%%%%%%%%%%%%%%%%%%%%%%%%%%%%%%%%%%%%%%%%%%%%%%
\subsection{Linear State-Space Systems}
\label{subsec:linear-state-space-control}
%%%%%%%%%%%%%%%%%%%%%%%%%%%%%%%%%%%%%%%%%%%%%%%%%%%%%%%%%%%%

A continuous-time linear time-invariant system has form

\[
\boxed{
\dot{\mathbf{x}}
=
A\mathbf{x}
+
B\mathbf{u},
}
\]

\[
\boxed{
\mathbf{y}
=
C\mathbf{x}
+
D\mathbf{u}.
}
\]

The matrices are:

\[
A
=
\text{state matrix},
\]

\[
B
=
\text{input matrix},
\]

\[
C
=
\text{output matrix},
\]

and

\[
D
=
\text{direct-feedthrough matrix}.
\]

%%%%%%%%%%%%%%%%%%%%%%%%%%%%%%%%%%%%%%%%%%%%%%%%%%%%%%%%%%%%
\subsection{State Equation Solution}
\label{subsec:state-equation-solution}
%%%%%%%%%%%%%%%%%%%%%%%%%%%%%%%%%%%%%%%%%%%%%%%%%%%%%%%%%%%%

For

\[
\dot{\mathbf{x}}
=
A\mathbf{x}
+
B\mathbf{u}(t),
\]

the solution is

\[
\boxed{
\mathbf{x}(t)
=
e^{At}\mathbf{x}_0
+
\int_0^t
e^{A(t-\tau)}
B\mathbf{u}(\tau)
\,d\tau.
}
\]

The matrix exponential

\[
e^{At}
\]

plays the role that \(e^{at}\) plays in scalar differential equations.

%%%%%%%%%%%%%%%%%%%%%%%%%%%%%%%%%%%%%%%%%%%%%%%%%%%%%%%%%%%%
\subsection{Worked Example: Scalar Controlled System}
\label{subsec:worked-scalar-control}
%%%%%%%%%%%%%%%%%%%%%%%%%%%%%%%%%%%%%%%%%%%%%%%%%%%%%%%%%%%%

\begin{workedexamplebox}{Constant Control Input}

Consider

\[
\dot{x}
=
-2x+u
\]

with constant input

\[
u=4
\]

and

\[
x(0)=0.
\]

Then

\[
\dot{x}+2x=4.
\]

The equilibrium is

\[
x^*=2.
\]

The solution is

\[
x(t)
=
2+Ce^{-2t}.
\]

Using \(x(0)=0\),

\[
C=-2.
\]

Therefore,

\begin{answerbox}
\[
\boxed{
x(t)
=
2(1-e^{-2t}).
}
\]
\end{answerbox}

The state approaches \(2\).

\end{workedexamplebox}

%%%%%%%%%%%%%%%%%%%%%%%%%%%%%%%%%%%%%%%%%%%%%%%%%%%%%%%%%%%%
\subsection{Equilibrium Points}
\label{subsec:control-equilibrium-points}
%%%%%%%%%%%%%%%%%%%%%%%%%%%%%%%%%%%%%%%%%%%%%%%%%%%%%%%%%%%%

For

\[
\dot{\mathbf{x}}
=
f(\mathbf{x},\mathbf{u}),
\]

an equilibrium associated with constant input \(\mathbf{u}^*\) satisfies

\[
\boxed{
f(\mathbf{x}^*,\mathbf{u}^*)=0.
}
\]

Control often seeks to stabilize a selected equilibrium.

%%%%%%%%%%%%%%%%%%%%%%%%%%%%%%%%%%%%%%%%%%%%%%%%%%%%%%%%%%%%
\subsection{Linearization}
\label{subsec:control-linearization}
%%%%%%%%%%%%%%%%%%%%%%%%%%%%%%%%%%%%%%%%%%%%%%%%%%%%%%%%%%%%

Suppose

\[
\dot{\mathbf{x}}
=
f(\mathbf{x},\mathbf{u}).
\]

Near equilibrium

\[
(\mathbf{x}^*,\mathbf{u}^*),
\]

define perturbations

\[
\delta\mathbf{x}
=
\mathbf{x}-\mathbf{x}^*,
\]

\[
\delta\mathbf{u}
=
\mathbf{u}-\mathbf{u}^*.
\]

The linearized system is

\[
\boxed{
\delta\dot{\mathbf{x}}
=
A\delta\mathbf{x}
+
B\delta\mathbf{u},
}
\]

where

\[
\boxed{
A
=
\left.
\frac{\partial f}{\partial\mathbf{x}}
\right|_{(\mathbf{x}^*,\mathbf{u}^*)},
}
\]

and

\[
\boxed{
B
=
\left.
\frac{\partial f}{\partial\mathbf{u}}
\right|_{(\mathbf{x}^*,\mathbf{u}^*)}.
}
\]

%%%%%%%%%%%%%%%%%%%%%%%%%%%%%%%%%%%%%%%%%%%%%%%%%%%%%%%%%%%%
\subsection{Transfer Functions}
\label{subsec:control-transfer-functions}
%%%%%%%%%%%%%%%%%%%%%%%%%%%%%%%%%%%%%%%%%%%%%%%%%%%%%%%%%%%%

For a linear time-invariant system with zero initial conditions, the transfer
function relates Laplace transforms of output and input:

\[
\boxed{
G(s)
=
\frac{Y(s)}{U(s)}.
}
\]

Transfer functions provide an input-output description without explicitly
displaying internal state variables.

%%%%%%%%%%%%%%%%%%%%%%%%%%%%%%%%%%%%%%%%%%%%%%%%%%%%%%%%%%%%
\subsection{Transfer Function From a Differential Equation}
\label{subsec:transfer-function-differential-equation}
%%%%%%%%%%%%%%%%%%%%%%%%%%%%%%%%%%%%%%%%%%%%%%%%%%%%%%%%%%%%

Consider

\[
a_2\ddot{y}
+
a_1\dot{y}
+
a_0y
=
b_1\dot{u}
+
b_0u.
\]

With zero initial conditions,

\[
(a_2s^2+a_1s+a_0)Y(s)
=
(b_1s+b_0)U(s).
\]

Therefore,

\[
\boxed{
G(s)
=
\frac{b_1s+b_0}
{a_2s^2+a_1s+a_0}.
}
\]

%%%%%%%%%%%%%%%%%%%%%%%%%%%%%%%%%%%%%%%%%%%%%%%%%%%%%%%%%%%%
\subsection{Transfer Function From State Space}
\label{subsec:state-space-transfer-function}
%%%%%%%%%%%%%%%%%%%%%%%%%%%%%%%%%%%%%%%%%%%%%%%%%%%%%%%%%%%%

For

\[
\dot{\mathbf{x}}
=
A\mathbf{x}+B\mathbf{u},
\]

\[
\mathbf{y}
=
C\mathbf{x}+D\mathbf{u},
\]

the transfer matrix is

\[
\boxed{
G(s)
=
C(sI-A)^{-1}B+D.
}
\]

This connects internal state dynamics to input-output behavior.

%%%%%%%%%%%%%%%%%%%%%%%%%%%%%%%%%%%%%%%%%%%%%%%%%%%%%%%%%%%%
\subsection{Poles}
\label{subsec:control-poles}
%%%%%%%%%%%%%%%%%%%%%%%%%%%%%%%%%%%%%%%%%%%%%%%%%%%%%%%%%%%%

Poles are values of \(s\) at which the transfer function becomes singular.

For a minimal state-space realization, poles correspond to eigenvalues of
\(A\).

The characteristic equation is

\[
\boxed{
\det(sI-A)=0.
}
\]

Pole locations strongly determine temporal behavior.

%%%%%%%%%%%%%%%%%%%%%%%%%%%%%%%%%%%%%%%%%%%%%%%%%%%%%%%%%%%%
\subsection{Zeros}
\label{subsec:control-zeros}
%%%%%%%%%%%%%%%%%%%%%%%%%%%%%%%%%%%%%%%%%%%%%%%%%%%%%%%%%%%%

Zeros are values of \(s\) that cause the transfer function numerator to
vanish.

For scalar

\[
G(s)
=
\frac{N(s)}{D(s)},
\]

zeros satisfy

\[
\boxed{
N(s)=0.
}
\]

Zeros affect transient response, frequency response, and achievable control
performance.

%%%%%%%%%%%%%%%%%%%%%%%%%%%%%%%%%%%%%%%%%%%%%%%%%%%%%%%%%%%%
\subsection{Continuous-Time Linear Stability}
\label{subsec:continuous-linear-stability}
%%%%%%%%%%%%%%%%%%%%%%%%%%%%%%%%%%%%%%%%%%%%%%%%%%%%%%%%%%%%

For the autonomous system

\[
\dot{\mathbf{x}}
=
A\mathbf{x},
\]

the equilibrium

\[
\mathbf{x}=0
\]

is asymptotically stable if every eigenvalue of \(A\) satisfies

\[
\boxed{
\operatorname{Re}(\lambda_i)<0.
}
\]

Such a matrix is called Hurwitz.

%%%%%%%%%%%%%%%%%%%%%%%%%%%%%%%%%%%%%%%%%%%%%%%%%%%%%%%%%%%%
\subsection{Worked Example: Eigenvalue Stability}
\label{subsec:worked-eigenvalue-stability}
%%%%%%%%%%%%%%%%%%%%%%%%%%%%%%%%%%%%%%%%%%%%%%%%%%%%%%%%%%%%

\begin{workedexamplebox}{Stable Linear System}

Let

\[
A
=
\begin{bmatrix}
-2 & 0\\
0 & -5
\end{bmatrix}.
\]

The eigenvalues are

\[
\lambda_1=-2,
\qquad
\lambda_2=-5.
\]

Both have negative real part.

\begin{answerbox}
\[
\boxed{
\mathbf{x}=0
\text{ is asymptotically stable}.
}
\]
\end{answerbox}

\end{workedexamplebox}

%%%%%%%%%%%%%%%%%%%%%%%%%%%%%%%%%%%%%%%%%%%%%%%%%%%%%%%%%%%%
\subsection{Unstable Linear Systems}
\label{subsec:unstable-linear-systems}
%%%%%%%%%%%%%%%%%%%%%%%%%%%%%%%%%%%%%%%%%%%%%%%%%%%%%%%%%%%%

If \(A\) has an eigenvalue satisfying

\[
\boxed{
\operatorname{Re}(\lambda)>0,
}
\]

then the equilibrium is unstable.

Solutions generally contain exponentially growing modes.

%%%%%%%%%%%%%%%%%%%%%%%%%%%%%%%%%%%%%%%%%%%%%%%%%%%%%%%%%%%%
\subsection{Marginal Cases}
\label{subsec:marginal-stability-control}
%%%%%%%%%%%%%%%%%%%%%%%%%%%%%%%%%%%%%%%%%%%%%%%%%%%%%%%%%%%%

Eigenvalues on the imaginary axis require additional care.

For example,

\[
\lambda=\pm i\omega
\]

can produce persistent oscillation.

Repeated or defective imaginary-axis eigenvalues can create unbounded
responses.

Thus zero real part alone does not guarantee stability.

%%%%%%%%%%%%%%%%%%%%%%%%%%%%%%%%%%%%%%%%%%%%%%%%%%%%%%%%%%%%
\subsection{BIBO Stability}
\label{subsec:bibo-stability}
%%%%%%%%%%%%%%%%%%%%%%%%%%%%%%%%%%%%%%%%%%%%%%%%%%%%%%%%%%%%

A system is bounded-input bounded-output stable if every bounded input
produces a bounded output.

For a proper rational continuous-time transfer function with no unstable
pole-zero cancellations in the realized input-output dynamics, BIBO
stability requires all relevant poles to lie strictly in the left half-plane.

%%%%%%%%%%%%%%%%%%%%%%%%%%%%%%%%%%%%%%%%%%%%%%%%%%%%%%%%%%%%
\subsection{Feedback State Control}
\label{subsec:state-feedback-control}
%%%%%%%%%%%%%%%%%%%%%%%%%%%%%%%%%%%%%%%%%%%%%%%%%%%%%%%%%%%%

Suppose

\[
\dot{\mathbf{x}}
=
A\mathbf{x}
+
B\mathbf{u}.
\]

Choose state feedback

\[
\boxed{
\mathbf{u}
=
-K\mathbf{x}.
}
\]

Then

\[
\boxed{
\dot{\mathbf{x}}
=
(A-BK)\mathbf{x}.
}
\]

The controller modifies the closed-loop eigenvalues.

%%%%%%%%%%%%%%%%%%%%%%%%%%%%%%%%%%%%%%%%%%%%%%%%%%%%%%%%%%%%
\subsection{Reference Tracking With State Feedback}
\label{subsec:state-feedback-reference}
%%%%%%%%%%%%%%%%%%%%%%%%%%%%%%%%%%%%%%%%%%%%%%%%%%%%%%%%%%%%

A simple tracking law may be written

\[
\boxed{
\mathbf{u}
=
-K\mathbf{x}
+
N\mathbf{r},
}
\]

where \(N\) is chosen to obtain an appropriate steady-state response to the
reference \(\mathbf{r}\).

State feedback primarily shapes the dynamics, while feedforward terms can
help set the desired equilibrium.

%%%%%%%%%%%%%%%%%%%%%%%%%%%%%%%%%%%%%%%%%%%%%%%%%%%%%%%%%%%%
\subsection{Pole Placement}
\label{subsec:pole-placement}
%%%%%%%%%%%%%%%%%%%%%%%%%%%%%%%%%%%%%%%%%%%%%%%%%%%%%%%%%%%%

If a system is controllable, one can assign the eigenvalues of

\[
A-BK
\]

to desired locations within important limitations of the model.

The goal is to choose

\[
\boxed{
K
}
\]

so that the closed-loop characteristic polynomial has specified roots.

%%%%%%%%%%%%%%%%%%%%%%%%%%%%%%%%%%%%%%%%%%%%%%%%%%%%%%%%%%%%
\subsection{Controllability}
\label{subsec:controllability}
%%%%%%%%%%%%%%%%%%%%%%%%%%%%%%%%%%%%%%%%%%%%%%%%%%%%%%%%%%%%

A state is controllable if an appropriate input can move the system to that
state in finite time.

For the LTI system

\[
\dot{\mathbf{x}}
=
A\mathbf{x}
+
B\mathbf{u},
\]

the controllability matrix is

\[
\boxed{
\mathcal{C}
=
\begin{bmatrix}
B & AB & A^2B & \cdots & A^{n-1}B
\end{bmatrix}.
}
\]

The system is controllable when

\[
\boxed{
\operatorname{rank}\mathcal{C}=n.
}
\]

%%%%%%%%%%%%%%%%%%%%%%%%%%%%%%%%%%%%%%%%%%%%%%%%%%%%%%%%%%%%
\subsection{Worked Example: Controllability}
\label{subsec:worked-controllability}
%%%%%%%%%%%%%%%%%%%%%%%%%%%%%%%%%%%%%%%%%%%%%%%%%%%%%%%%%%%%

\begin{workedexamplebox}{A Two-State System}

Let

\[
A
=
\begin{bmatrix}
0 & 1\\
0 & 0
\end{bmatrix},
\qquad
B
=
\begin{bmatrix}
0\\
1
\end{bmatrix}.
\]

Then

\[
AB
=
\begin{bmatrix}
1\\
0
\end{bmatrix}.
\]

Thus

\[
\mathcal{C}
=
\begin{bmatrix}
0 & 1\\
1 & 0
\end{bmatrix}.
\]

Since

\[
\operatorname{rank}\mathcal{C}=2,
\]

\begin{answerbox}
\[
\boxed{
\text{the system is controllable}.
}
\]
\end{answerbox}

\end{workedexamplebox}

%%%%%%%%%%%%%%%%%%%%%%%%%%%%%%%%%%%%%%%%%%%%%%%%%%%%%%%%%%%%
\subsection{Uncontrollable Modes}
\label{subsec:uncontrollable-modes}
%%%%%%%%%%%%%%%%%%%%%%%%%%%%%%%%%%%%%%%%%%%%%%%%%%%%%%%%%%%%

If

\[
\operatorname{rank}\mathcal{C}<n,
\]

some state directions cannot be influenced by the input.

An unstable uncontrollable mode cannot be stabilized through the available
control channel.

This illustrates why actuator placement matters.

%%%%%%%%%%%%%%%%%%%%%%%%%%%%%%%%%%%%%%%%%%%%%%%%%%%%%%%%%%%%
\subsection{Observability}
\label{subsec:observability}
%%%%%%%%%%%%%%%%%%%%%%%%%%%%%%%%%%%%%%%%%%%%%%%%%%%%%%%%%%%%

A system is observable if its initial state can be uniquely reconstructed
from input-output measurements over a finite interval.

For

\[
\dot{\mathbf{x}}
=
A\mathbf{x}
+
B\mathbf{u},
\]

\[
\mathbf{y}
=
C\mathbf{x},
\]

the observability matrix is

\[
\boxed{
\mathcal{O}
=
\begin{bmatrix}
C\\
CA\\
CA^2\\
\vdots\\
CA^{n-1}
\end{bmatrix}.
}
\]

The system is observable if

\[
\boxed{
\operatorname{rank}\mathcal{O}=n.
}
\]

%%%%%%%%%%%%%%%%%%%%%%%%%%%%%%%%%%%%%%%%%%%%%%%%%%%%%%%%%%%%
\subsection{Worked Example: Observability}
\label{subsec:worked-observability}
%%%%%%%%%%%%%%%%%%%%%%%%%%%%%%%%%%%%%%%%%%%%%%%%%%%%%%%%%%%%

\begin{workedexamplebox}{Observing Both States}

Let

\[
A
=
\begin{bmatrix}
0 & 1\\
0 & 0
\end{bmatrix},
\qquad
C
=
\begin{bmatrix}
1 & 0
\end{bmatrix}.
\]

Then

\[
CA
=
\begin{bmatrix}
0 & 1
\end{bmatrix}.
\]

Therefore,

\[
\mathcal{O}
=
\begin{bmatrix}
1 & 0\\
0 & 1
\end{bmatrix}.
\]

Hence,

\begin{answerbox}
\[
\boxed{
\operatorname{rank}\mathcal{O}=2,
}
\]

so the system is observable.
\end{answerbox}

\end{workedexamplebox}

%%%%%%%%%%%%%%%%%%%%%%%%%%%%%%%%%%%%%%%%%%%%%%%%%%%%%%%%%%%%
\subsection{State Observers}
\label{subsec:state-observers}
%%%%%%%%%%%%%%%%%%%%%%%%%%%%%%%%%%%%%%%%%%%%%%%%%%%%%%%%%%%%

Often not every state can be measured directly.

An observer estimates the state using known inputs and measured outputs.

A Luenberger observer has form

\[
\boxed{
\dot{\hat{\mathbf{x}}}
=
A\hat{\mathbf{x}}
+
B\mathbf{u}
+
L
\left(
\mathbf{y}
-
C\hat{\mathbf{x}}
\right).
}
\]

%%%%%%%%%%%%%%%%%%%%%%%%%%%%%%%%%%%%%%%%%%%%%%%%%%%%%%%%%%%%
\subsection{Observer Error Dynamics}
\label{subsec:observer-error-dynamics}
%%%%%%%%%%%%%%%%%%%%%%%%%%%%%%%%%%%%%%%%%%%%%%%%%%%%%%%%%%%%

Define estimation error

\[
\boxed{
\mathbf{e}
=
\mathbf{x}
-
\hat{\mathbf{x}}.
}
\]

Then

\[
\boxed{
\dot{\mathbf{e}}
=
(A-LC)\mathbf{e}.
}
\]

If the pair \((A,C)\) is observable, \(L\) can be selected to place observer
error poles.

%%%%%%%%%%%%%%%%%%%%%%%%%%%%%%%%%%%%%%%%%%%%%%%%%%%%%%%%%%%%
\subsection{Controller-Observer Combination}
\label{subsec:controller-observer-combination}
%%%%%%%%%%%%%%%%%%%%%%%%%%%%%%%%%%%%%%%%%%%%%%%%%%%%%%%%%%%%

When states are not directly measured, use estimated-state feedback:

\[
\boxed{
\mathbf{u}
=
-K\hat{\mathbf{x}}.
}
\]

For linear systems under standard controllability and observability
conditions, controller and observer dynamics can be designed separately.

This is related to the separation principle.

%%%%%%%%%%%%%%%%%%%%%%%%%%%%%%%%%%%%%%%%%%%%%%%%%%%%%%%%%%%%
\subsection{Proportional Control}
\label{subsec:proportional-control}
%%%%%%%%%%%%%%%%%%%%%%%%%%%%%%%%%%%%%%%%%%%%%%%%%%%%%%%%%%%%

A proportional controller is

\[
\boxed{
u(t)
=
K_Pe(t).
}
\]

Large \(K_P\) can reduce error and increase responsiveness.

However, excessive gain may cause overshoot, oscillation, noise
amplification, or instability.

%%%%%%%%%%%%%%%%%%%%%%%%%%%%%%%%%%%%%%%%%%%%%%%%%%%%%%%%%%%%
\subsection{Integral Control}
\label{subsec:integral-control}
%%%%%%%%%%%%%%%%%%%%%%%%%%%%%%%%%%%%%%%%%%%%%%%%%%%%%%%%%%%%

Integral control uses accumulated error:

\[
\boxed{
u(t)
=
K_I
\int_0^t
e(\tau)\,d\tau.
}
\]

Integral action is particularly useful for removing persistent steady-state
error.

%%%%%%%%%%%%%%%%%%%%%%%%%%%%%%%%%%%%%%%%%%%%%%%%%%%%%%%%%%%%
\subsection{Derivative Control}
\label{subsec:derivative-control}
%%%%%%%%%%%%%%%%%%%%%%%%%%%%%%%%%%%%%%%%%%%%%%%%%%%%%%%%%%%%

Derivative control responds to the rate of error change:

\[
\boxed{
u(t)
=
K_D
\frac{de}{dt}.
}
\]

Derivative action can improve damping but is sensitive to measurement noise.

Practical implementations typically use filtered derivatives.

%%%%%%%%%%%%%%%%%%%%%%%%%%%%%%%%%%%%%%%%%%%%%%%%%%%%%%%%%%%%
\subsection{PID Control}
\label{subsec:pid-control}
%%%%%%%%%%%%%%%%%%%%%%%%%%%%%%%%%%%%%%%%%%%%%%%%%%%%%%%%%%%%

A proportional-integral-derivative controller is

\[
\boxed{
u(t)
=
K_Pe(t)
+
K_I
\int_0^t
e(\tau)\,d\tau
+
K_D
\frac{de}{dt}.
}
\]

PID controllers are widely used because they combine simple structure with
effective performance in many practical systems.

%%%%%%%%%%%%%%%%%%%%%%%%%%%%%%%%%%%%%%%%%%%%%%%%%%%%%%%%%%%%
\subsection{Worked Example: Proportional Feedback}
\label{subsec:worked-proportional-feedback}
%%%%%%%%%%%%%%%%%%%%%%%%%%%%%%%%%%%%%%%%%%%%%%%%%%%%%%%%%%%%

\begin{workedexamplebox}{Stabilizing a Scalar System}

Consider

\[
\dot{x}
=
2x+u.
\]

Without control, the system is unstable.

Choose

\[
u=-Kx.
\]

Then

\[
\dot{x}
=
(2-K)x.
\]

For asymptotic stability,

\[
2-K<0.
\]

Therefore,

\begin{answerbox}
\[
\boxed{
K>2.
}
\]
\end{answerbox}

For example, \(K=5\) gives

\[
\dot{x}=-3x.
\]

\end{workedexamplebox}

%%%%%%%%%%%%%%%%%%%%%%%%%%%%%%%%%%%%%%%%%%%%%%%%%%%%%%%%%%%%
\subsection{Steady-State Error}
\label{subsec:steady-state-error}
%%%%%%%%%%%%%%%%%%%%%%%%%%%%%%%%%%%%%%%%%%%%%%%%%%%%%%%%%%%%

For tracking error

\[
e(t)=r(t)-y(t),
\]

the steady-state error is

\[
\boxed{
e_{\mathrm{ss}}
=
\lim_{t\to\infty}e(t),
}
\]

when the limit exists.

Control design often seeks

\[
\boxed{
e_{\mathrm{ss}}=0.
}
\]

Integral action is a common mechanism for achieving zero steady-state error
against certain constant references and disturbances.

%%%%%%%%%%%%%%%%%%%%%%%%%%%%%%%%%%%%%%%%%%%%%%%%%%%%%%%%%%%%
\subsection{Transient Performance}
\label{subsec:control-transient-performance}
%%%%%%%%%%%%%%%%%%%%%%%%%%%%%%%%%%%%%%%%%%%%%%%%%%%%%%%%%%%%

Important transient measures include:

\[
\boxed{
\text{rise time},
}
\]

\[
\boxed{
\text{peak time},
}
\]

\[
\boxed{
\text{overshoot},
}
\]

and

\[
\boxed{
\text{settling time}.
}
\]

These quantify how quickly and smoothly a controlled system approaches its
desired behavior.

%%%%%%%%%%%%%%%%%%%%%%%%%%%%%%%%%%%%%%%%%%%%%%%%%%%%%%%%%%%%
\subsection{Second-Order Standard Form}
\label{subsec:second-order-control-standard-form}
%%%%%%%%%%%%%%%%%%%%%%%%%%%%%%%%%%%%%%%%%%%%%%%%%%%%%%%%%%%%

A common second-order transfer function is

\[
\boxed{
G(s)
=
\frac{\omega_n^2}
{s^2+2\zeta\omega_ns+\omega_n^2}.
}
\]

Here:

\[
\omega_n
=
\text{natural frequency},
\]

and

\[
\zeta
=
\text{damping ratio}.
\]

These parameters strongly influence transient response.

%%%%%%%%%%%%%%%%%%%%%%%%%%%%%%%%%%%%%%%%%%%%%%%%%%%%%%%%%%%%
\subsection{Damping Ratio}
\label{subsec:control-damping-ratio}
%%%%%%%%%%%%%%%%%%%%%%%%%%%%%%%%%%%%%%%%%%%%%%%%%%%%%%%%%%%%

For

\[
0<\zeta<1,
\]

the system is underdamped.

For

\[
\zeta=1,
\]

it is critically damped.

For

\[
\zeta>1,
\]

it is overdamped.

Increasing damping generally reduces oscillation but may slow response.

%%%%%%%%%%%%%%%%%%%%%%%%%%%%%%%%%%%%%%%%%%%%%%%%%%%%%%%%%%%%
\subsection{Characteristic Equation}
\label{subsec:control-characteristic-equation}
%%%%%%%%%%%%%%%%%%%%%%%%%%%%%%%%%%%%%%%%%%%%%%%%%%%%%%%%%%%%

Closed-loop poles are roots of the characteristic equation.

For unity negative feedback with forward transfer function \(K G(s)\),

\[
\boxed{
1+KG(s)=0.
}
\]

Changing controller gain moves the closed-loop poles.

%%%%%%%%%%%%%%%%%%%%%%%%%%%%%%%%%%%%%%%%%%%%%%%%%%%%%%%%%%%%
\subsection{Root Locus}
\label{subsec:root-locus}
%%%%%%%%%%%%%%%%%%%%%%%%%%%%%%%%%%%%%%%%%%%%%%%%%%%%%%%%%%%%

A root locus shows how closed-loop poles move as a controller gain varies.

It provides geometric insight into:

\[
\boxed{
\text{stability},
}
\]

\[
\boxed{
\text{damping},
}
\]

and

\[
\boxed{
\text{transient response}.
}
\]

Root-locus methods are especially useful for classical single-input
single-output design.

%%%%%%%%%%%%%%%%%%%%%%%%%%%%%%%%%%%%%%%%%%%%%%%%%%%%%%%%%%%%
\subsection{Frequency Response}
\label{subsec:control-frequency-response}
%%%%%%%%%%%%%%%%%%%%%%%%%%%%%%%%%%%%%%%%%%%%%%%%%%%%%%%%%%%%

For transfer function

\[
G(s),
\]

evaluate

\[
\boxed{
G(i\omega).
}
\]

Its magnitude and phase describe how sinusoidal inputs are amplified and
shifted at frequency \(\omega\).

Frequency-domain methods are central in robust feedback design.

%%%%%%%%%%%%%%%%%%%%%%%%%%%%%%%%%%%%%%%%%%%%%%%%%%%%%%%%%%%%
\subsection{Bode Magnitude and Phase}
\label{subsec:bode-plots}
%%%%%%%%%%%%%%%%%%%%%%%%%%%%%%%%%%%%%%%%%%%%%%%%%%%%%%%%%%%%

A Bode plot displays

\[
\boxed{
20\log_{10}|G(i\omega)|
}
\]

against logarithmic frequency, together with

\[
\boxed{
\arg G(i\omega).
}
\]

This representation makes products of transfer functions easier to analyze
because magnitudes add on a decibel scale.

%%%%%%%%%%%%%%%%%%%%%%%%%%%%%%%%%%%%%%%%%%%%%%%%%%%%%%%%%%%%
\subsection{Gain Margin}
\label{subsec:gain-margin}
%%%%%%%%%%%%%%%%%%%%%%%%%%%%%%%%%%%%%%%%%%%%%%%%%%%%%%%%%%%%

Gain margin measures how much loop gain can increase before a nominal
closed-loop system reaches an instability boundary, under the standard
frequency-response interpretation.

A larger positive gain margin typically indicates greater tolerance to gain
uncertainty.

%%%%%%%%%%%%%%%%%%%%%%%%%%%%%%%%%%%%%%%%%%%%%%%%%%%%%%%%%%%%
\subsection{Phase Margin}
\label{subsec:phase-margin}
%%%%%%%%%%%%%%%%%%%%%%%%%%%%%%%%%%%%%%%%%%%%%%%%%%%%%%%%%%%%

Phase margin measures additional phase lag that can be introduced before the
feedback system reaches the instability boundary.

It is often used as a practical robustness indicator.

%%%%%%%%%%%%%%%%%%%%%%%%%%%%%%%%%%%%%%%%%%%%%%%%%%%%%%%%%%%%
\subsection{Feedforward Control}
\label{subsec:feedforward-control}
%%%%%%%%%%%%%%%%%%%%%%%%%%%%%%%%%%%%%%%%%%%%%%%%%%%%%%%%%%%%

Feedforward control uses known or measured information about desired motion
or disturbances before output error develops.

Schematically,

\[
\boxed{
u
=
u_{\mathrm{ff}}
+
u_{\mathrm{fb}}.
}
\]

Feedforward improves nominal response.

Feedback corrects uncertainty and disturbances.

%%%%%%%%%%%%%%%%%%%%%%%%%%%%%%%%%%%%%%%%%%%%%%%%%%%%%%%%%%%%
\subsection{Disturbance Rejection}
\label{subsec:disturbance-rejection}
%%%%%%%%%%%%%%%%%%%%%%%%%%%%%%%%%%%%%%%%%%%%%%%%%%%%%%%%%%%%

Suppose a disturbance \(d(t)\) affects the plant.

A feedback controller should reduce its influence on the output.

The closed-loop mapping from disturbance to output is therefore an important
design quantity.

%%%%%%%%%%%%%%%%%%%%%%%%%%%%%%%%%%%%%%%%%%%%%%%%%%%%%%%%%%%%
\subsection{Sensitivity Function}
\label{subsec:control-sensitivity-function}
%%%%%%%%%%%%%%%%%%%%%%%%%%%%%%%%%%%%%%%%%%%%%%%%%%%%%%%%%%%%

For scalar loop transfer function

\[
L(s)=C(s)P(s),
\]

the sensitivity function is

\[
\boxed{
S(s)
=
\frac{1}{1+L(s)}.
}
\]

Small \(|S(i\omega)|\) indicates reduced sensitivity to certain disturbances
and modeling errors at frequency \(\omega\).

%%%%%%%%%%%%%%%%%%%%%%%%%%%%%%%%%%%%%%%%%%%%%%%%%%%%%%%%%%%%
\subsection{Complementary Sensitivity}
\label{subsec:complementary-sensitivity}
%%%%%%%%%%%%%%%%%%%%%%%%%%%%%%%%%%%%%%%%%%%%%%%%%%%%%%%%%%%%

The complementary sensitivity function is

\[
\boxed{
T(s)
=
\frac{L(s)}
{1+L(s)}.
}
\]

These satisfy

\[
\boxed{
S(s)+T(s)=1.
}
\]

Control design balances disturbance rejection, tracking, uncertainty, and
noise amplification across frequency.

%%%%%%%%%%%%%%%%%%%%%%%%%%%%%%%%%%%%%%%%%%%%%%%%%%%%%%%%%%%%
\subsection{Lyapunov Stability}
\label{subsec:control-lyapunov-stability}
%%%%%%%%%%%%%%%%%%%%%%%%%%%%%%%%%%%%%%%%%%%%%%%%%%%%%%%%%%%%

For

\[
\dot{\mathbf{x}}
=
f(\mathbf{x}),
\]

a Lyapunov function \(V(\mathbf{x})\) is a scalar function used to measure
distance-like behavior relative to an equilibrium.

Typical conditions are

\[
\boxed{
V(\mathbf{x})>0
\quad
\text{for }\mathbf{x}\neq0,
}
\]

\[
\boxed{
V(0)=0,
}
\]

and

\[
\boxed{
\dot V(\mathbf{x})<0
\quad
\text{for }\mathbf{x}\neq0.
}
\]

These conditions imply asymptotic stability under standard hypotheses.

%%%%%%%%%%%%%%%%%%%%%%%%%%%%%%%%%%%%%%%%%%%%%%%%%%%%%%%%%%%%
\subsection{Worked Example: Lyapunov Function}
\label{subsec:worked-lyapunov-control}
%%%%%%%%%%%%%%%%%%%%%%%%%%%%%%%%%%%%%%%%%%%%%%%%%%%%%%%%%%%%

\begin{workedexamplebox}{Scalar Nonlinear Stability}

Consider

\[
\dot{x}
=
-x^3.
\]

Choose

\[
V(x)
=
\frac12x^2.
\]

Then

\[
V(x)>0
\]

for \(x\neq0\), and

\[
V(0)=0.
\]

Also,

\[
\dot V
=
x\dot{x}
=
-x^4.
\]

Thus

\[
\dot V<0
\]

for \(x\neq0\).

\begin{answerbox}
\[
\boxed{
x=0
\text{ is asymptotically stable}.
}
\]
\end{answerbox}

\end{workedexamplebox}

%%%%%%%%%%%%%%%%%%%%%%%%%%%%%%%%%%%%%%%%%%%%%%%%%%%%%%%%%%%%
\subsection{Lyapunov Equation for Linear Systems}
\label{subsec:linear-lyapunov-equation}
%%%%%%%%%%%%%%%%%%%%%%%%%%%%%%%%%%%%%%%%%%%%%%%%%%%%%%%%%%%%

For

\[
\dot{\mathbf{x}}
=
A\mathbf{x},
\]

seek

\[
\boxed{
V(\mathbf{x})
=
\mathbf{x}^{T}P\mathbf{x},
}
\]

with \(P\) positive definite.

Then

\[
\dot V
=
\mathbf{x}^{T}
(A^TP+PA)
\mathbf{x}.
\]

Choose \(P\) satisfying

\[
\boxed{
A^TP+PA
=
-Q,
}
\]

where \(Q\) is positive definite.

For Hurwitz \(A\), this equation has a unique positive-definite solution
\(P\) for each positive-definite \(Q\).

%%%%%%%%%%%%%%%%%%%%%%%%%%%%%%%%%%%%%%%%%%%%%%%%%%%%%%%%%%%%
\subsection{Optimal Control}
\label{subsec:control-optimal-control}
%%%%%%%%%%%%%%%%%%%%%%%%%%%%%%%%%%%%%%%%%%%%%%%%%%%%%%%%%%%%

Optimal control chooses a time-dependent input to minimize a performance
criterion subject to system dynamics.

A general problem is

\[
\boxed{
\min_{\mathbf{u}(\cdot)}
\left[
\Phi(\mathbf{x}(T))
+
\int_0^T
L(\mathbf{x},\mathbf{u},t)
\,dt
\right]
}
\]

subject to

\[
\boxed{
\dot{\mathbf{x}}
=
f(\mathbf{x},\mathbf{u},t).
}
\]

%%%%%%%%%%%%%%%%%%%%%%%%%%%%%%%%%%%%%%%%%%%%%%%%%%%%%%%%%%%%
\subsection{Quadratic Performance Index}
\label{subsec:quadratic-performance-index}
%%%%%%%%%%%%%%%%%%%%%%%%%%%%%%%%%%%%%%%%%%%%%%%%%%%%%%%%%%%%

For linear systems, a common objective is

\[
\boxed{
J
=
\int_0^\infty
\left(
\mathbf{x}^TQ\mathbf{x}
+
\mathbf{u}^TR\mathbf{u}
\right)
dt.
}
\]

The matrix \(Q\) penalizes state deviation.

The matrix \(R\) penalizes control effort.

%%%%%%%%%%%%%%%%%%%%%%%%%%%%%%%%%%%%%%%%%%%%%%%%%%%%%%%%%%%%
\subsection{Linear-Quadratic Regulator}
\label{subsec:linear-quadratic-regulator}
%%%%%%%%%%%%%%%%%%%%%%%%%%%%%%%%%%%%%%%%%%%%%%%%%%%%%%%%%%%%

Consider

\[
\dot{\mathbf{x}}
=
A\mathbf{x}
+
B\mathbf{u}
\]

and quadratic infinite-horizon cost.

Under standard stabilizability and detectability conditions, the optimal
control has form

\[
\boxed{
\mathbf{u}
=
-K\mathbf{x},
}
\]

where

\[
\boxed{
K
=
R^{-1}B^TP.
}
\]

The matrix \(P\) satisfies an algebraic Riccati equation.

%%%%%%%%%%%%%%%%%%%%%%%%%%%%%%%%%%%%%%%%%%%%%%%%%%%%%%%%%%%%
\subsection{Algebraic Riccati Equation}
\label{subsec:algebraic-riccati-equation}
%%%%%%%%%%%%%%%%%%%%%%%%%%%%%%%%%%%%%%%%%%%%%%%%%%%%%%%%%%%%

The continuous-time algebraic Riccati equation is

\[
\boxed{
A^TP
+
PA
-
PBR^{-1}B^TP
+
Q
=
0.
}
\]

Its stabilizing solution determines the LQR feedback gain.

%%%%%%%%%%%%%%%%%%%%%%%%%%%%%%%%%%%%%%%%%%%%%%%%%%%%%%%%%%%%
\subsection{Worked Example: Scalar LQR}
\label{subsec:worked-scalar-lqr}
%%%%%%%%%%%%%%%%%%%%%%%%%%%%%%%%%%%%%%%%%%%%%%%%%%%%%%%%%%%%

\begin{workedexamplebox}{Optimal Feedback for a Scalar Integrator}

Consider

\[
\dot{x}=u
\]

with cost

\[
J
=
\int_0^\infty
(x^2+u^2)\,dt.
\]

Here,

\[
A=0,
\qquad
B=1,
\qquad
Q=1,
\qquad
R=1.
\]

The Riccati equation becomes

\[
-P^2+1=0.
\]

The positive stabilizing solution is

\[
P=1.
\]

Thus

\[
K=1.
\]

Therefore,

\begin{answerbox}
\[
\boxed{
u=-x.
}
\]
\end{answerbox}

The closed-loop system is

\[
\dot{x}=-x.
\]

\end{workedexamplebox}

%%%%%%%%%%%%%%%%%%%%%%%%%%%%%%%%%%%%%%%%%%%%%%%%%%%%%%%%%%%%
\subsection{Finite-Horizon LQR}
\label{subsec:finite-horizon-lqr}
%%%%%%%%%%%%%%%%%%%%%%%%%%%%%%%%%%%%%%%%%%%%%%%%%%%%%%%%%%%%

For finite time horizon \(T\), the optimal feedback generally becomes
time-dependent:

\[
\boxed{
\mathbf{u}(t)
=
-R^{-1}B^TP(t)\mathbf{x}(t).
}
\]

The matrix \(P(t)\) satisfies a differential Riccati equation solved
backward from a terminal condition.

%%%%%%%%%%%%%%%%%%%%%%%%%%%%%%%%%%%%%%%%%%%%%%%%%%%%%%%%%%%%
\subsection{Pontryagin Maximum Principle}
\label{subsec:pontryagin-maximum-principle}
%%%%%%%%%%%%%%%%%%%%%%%%%%%%%%%%%%%%%%%%%%%%%%%%%%%%%%%%%%%%

For system

\[
\dot{\mathbf{x}}
=
f(\mathbf{x},\mathbf{u},t),
\]

define a control Hamiltonian

\[
\boxed{
\mathcal{H}
=
L(\mathbf{x},\mathbf{u},t)
+
\boldsymbol{\lambda}^{T}
f(\mathbf{x},\mathbf{u},t).
}
\]

Necessary conditions for an optimum include state dynamics, costate
dynamics, and an optimality condition for \(\mathbf{u}\).

%%%%%%%%%%%%%%%%%%%%%%%%%%%%%%%%%%%%%%%%%%%%%%%%%%%%%%%%%%%%
\subsection{Costate Equation}
\label{subsec:costate-equation}
%%%%%%%%%%%%%%%%%%%%%%%%%%%%%%%%%%%%%%%%%%%%%%%%%%%%%%%%%%%%

The costate satisfies

\[
\boxed{
\dot{\boldsymbol{\lambda}}
=
-
\frac{\partial\mathcal{H}}
{\partial\mathbf{x}}.
}
\]

The state satisfies

\[
\boxed{
\dot{\mathbf{x}}
=
\frac{\partial\mathcal{H}}
{\partial\boldsymbol{\lambda}}.
}
\]

This creates a coupled two-point boundary-value problem in many optimal
control problems.

%%%%%%%%%%%%%%%%%%%%%%%%%%%%%%%%%%%%%%%%%%%%%%%%%%%%%%%%%%%%
\subsection{Hamilton--Jacobi--Bellman Equation}
\label{subsec:hjb-control}
%%%%%%%%%%%%%%%%%%%%%%%%%%%%%%%%%%%%%%%%%%%%%%%%%%%%%%%%%%%%

Dynamic programming describes optimal control through a value function

\[
V(\mathbf{x},t).
\]

The continuous-time Hamilton--Jacobi--Bellman equation has schematic form

\[
\boxed{
-V_t
=
\min_{\mathbf{u}}
\left\{
L(\mathbf{x},\mathbf{u},t)
+
\nabla V
\cdot
f(\mathbf{x},\mathbf{u},t)
\right\}.
}
\]

The HJB equation characterizes optimal feedback directly.

%%%%%%%%%%%%%%%%%%%%%%%%%%%%%%%%%%%%%%%%%%%%%%%%%%%%%%%%%%%%
\subsection{Model Predictive Control}
\label{subsec:model-predictive-control}
%%%%%%%%%%%%%%%%%%%%%%%%%%%%%%%%%%%%%%%%%%%%%%%%%%%%%%%%%%%%

Model predictive control repeatedly solves a finite-horizon optimization
problem.

At each decision time:

\begin{enumerate}
    \item measure or estimate the current state;
    \item predict future system behavior;
    \item optimize a sequence of future controls;
    \item apply only the first control action;
    \item shift the horizon forward;
    \item repeat.
\end{enumerate}

This is also called receding-horizon control.

%%%%%%%%%%%%%%%%%%%%%%%%%%%%%%%%%%%%%%%%%%%%%%%%%%%%%%%%%%%%
\subsection{Why Model Predictive Control Is Useful}
\label{subsec:mpc-advantages}
%%%%%%%%%%%%%%%%%%%%%%%%%%%%%%%%%%%%%%%%%%%%%%%%%%%%%%%%%%%%

MPC can explicitly handle constraints such as

\[
\boxed{
u_{\min}
\leq
u_k
\leq
u_{\max},
}
\]

and

\[
\boxed{
x_{\min}
\leq
x_k
\leq
x_{\max}.
}
\]

This makes it attractive for industrial processes, vehicles, energy systems,
and constrained multivariable control.

%%%%%%%%%%%%%%%%%%%%%%%%%%%%%%%%%%%%%%%%%%%%%%%%%%%%%%%%%%%%
\subsection{Control Constraints}
\label{subsec:control-constraints}
%%%%%%%%%%%%%%%%%%%%%%%%%%%%%%%%%%%%%%%%%%%%%%%%%%%%%%%%%%%%

Real actuators have limits.

Examples include:

\[
\boxed{
\text{maximum force},
}
\]

\[
\boxed{
\text{maximum voltage},
}
\]

\[
\boxed{
\text{maximum steering angle},
}
\]

and

\[
\boxed{
\text{rate limits}.
}
\]

Ignoring constraints can make an otherwise mathematically valid controller
physically impossible to implement.

%%%%%%%%%%%%%%%%%%%%%%%%%%%%%%%%%%%%%%%%%%%%%%%%%%%%%%%%%%%%
\subsection{Actuator Saturation}
\label{subsec:actuator-saturation}
%%%%%%%%%%%%%%%%%%%%%%%%%%%%%%%%%%%%%%%%%%%%%%%%%%%%%%%%%%%%

A saturated actuator may satisfy

\[
\boxed{
u
=
\operatorname{sat}(u_c),
}
\]

where \(u_c\) is the requested control and

\[
\operatorname{sat}(u_c)
=
\begin{cases}
u_{\max}, & u_c>u_{\max},\\
u_c, & u_{\min}\leq u_c\leq u_{\max},\\
u_{\min}, & u_c<u_{\min}.
\end{cases}
\]

Saturation introduces nonlinearity.

%%%%%%%%%%%%%%%%%%%%%%%%%%%%%%%%%%%%%%%%%%%%%%%%%%%%%%%%%%%%
\subsection{Integral Windup}
\label{subsec:integral-windup}
%%%%%%%%%%%%%%%%%%%%%%%%%%%%%%%%%%%%%%%%%%%%%%%%%%%%%%%%%%%%

When an actuator saturates, an integral controller can continue accumulating
error.

This may produce a large internal integral state and poor recovery after the
system leaves saturation.

This phenomenon is called integral windup.

Anti-windup methods limit or correct the integral state.

%%%%%%%%%%%%%%%%%%%%%%%%%%%%%%%%%%%%%%%%%%%%%%%%%%%%%%%%%%%%
\subsection{Robust Control}
\label{subsec:robust-control}
%%%%%%%%%%%%%%%%%%%%%%%%%%%%%%%%%%%%%%%%%%%%%%%%%%%%%%%%%%%%

Real models contain uncertainty.

A robust controller seeks acceptable performance despite:

\[
\boxed{
\text{parameter variation},
}
\]

\[
\boxed{
\text{unmodeled dynamics},
}
\]

\[
\boxed{
\text{disturbances},
}
\]

and

\[
\boxed{
\text{measurement noise}.
}
\]

Stability under the nominal model alone is not always sufficient.

%%%%%%%%%%%%%%%%%%%%%%%%%%%%%%%%%%%%%%%%%%%%%%%%%%%%%%%%%%%%
\subsection{Parametric Uncertainty}
\label{subsec:parametric-uncertainty-control}
%%%%%%%%%%%%%%%%%%%%%%%%%%%%%%%%%%%%%%%%%%%%%%%%%%%%%%%%%%%%

Suppose a model depends on uncertain parameter \(\theta\):

\[
\boxed{
\dot{\mathbf{x}}
=
A(\theta)\mathbf{x}
+
B(\theta)\mathbf{u}.
}
\]

A robust design may seek stability for all

\[
\boxed{
\theta
\in
\Theta.
}
\]

This converts control design into a family of stability problems.

%%%%%%%%%%%%%%%%%%%%%%%%%%%%%%%%%%%%%%%%%%%%%%%%%%%%%%%%%%%%
\subsection{Stochastic Control}
\label{subsec:stochastic-control}
%%%%%%%%%%%%%%%%%%%%%%%%%%%%%%%%%%%%%%%%%%%%%%%%%%%%%%%%%%%%

A stochastic system may satisfy

\[
\boxed{
d\mathbf{x}
=
f(\mathbf{x},\mathbf{u},t)\,dt
+
G(\mathbf{x},t)\,d\mathbf{W}_t.
}
\]

The controller must account for random disturbances.

Objectives commonly involve expected cost:

\[
\boxed{
\min
\mathbb{E}[J].
}
\]

%%%%%%%%%%%%%%%%%%%%%%%%%%%%%%%%%%%%%%%%%%%%%%%%%%%%%%%%%%%%
\subsection{State Estimation Under Noise}
\label{subsec:state-estimation-noise}
%%%%%%%%%%%%%%%%%%%%%%%%%%%%%%%%%%%%%%%%%%%%%%%%%%%%%%%%%%%%

Suppose

\[
\boxed{
\mathbf{x}_{k+1}
=
A\mathbf{x}_k
+
B\mathbf{u}_k
+
\mathbf{w}_k,
}
\]

and

\[
\boxed{
\mathbf{y}_k
=
C\mathbf{x}_k
+
\mathbf{v}_k,
}
\]

where \(\mathbf{w}_k\) and \(\mathbf{v}_k\) are process and measurement
noise.

One seeks an estimate

\[
\boxed{
\hat{\mathbf{x}}_k.
}
\]

%%%%%%%%%%%%%%%%%%%%%%%%%%%%%%%%%%%%%%%%%%%%%%%%%%%%%%%%%%%%
\subsection{Kalman Filter}
\label{subsec:kalman-filter}
%%%%%%%%%%%%%%%%%%%%%%%%%%%%%%%%%%%%%%%%%%%%%%%%%%%%%%%%%%%%

For linear systems with standard Gaussian noise assumptions, the Kalman
filter recursively combines:

\[
\boxed{
\text{model prediction}
}
\]

and

\[
\boxed{
\text{measurement correction}.
}
\]

A prediction step estimates the next state.

A correction step adjusts that estimate based on the measurement residual.

%%%%%%%%%%%%%%%%%%%%%%%%%%%%%%%%%%%%%%%%%%%%%%%%%%%%%%%%%%%%
\subsection{Kalman Measurement Update}
\label{subsec:kalman-measurement-update}
%%%%%%%%%%%%%%%%%%%%%%%%%%%%%%%%%%%%%%%%%%%%%%%%%%%%%%%%%%%%

A conceptual measurement update is

\[
\boxed{
\hat{\mathbf{x}}_{k|k}
=
\hat{\mathbf{x}}_{k|k-1}
+
K_k
\left(
\mathbf{y}_k
-
C\hat{\mathbf{x}}_{k|k-1}
\right).
}
\]

The matrix \(K_k\) is the Kalman gain.

It balances model uncertainty against measurement uncertainty.

%%%%%%%%%%%%%%%%%%%%%%%%%%%%%%%%%%%%%%%%%%%%%%%%%%%%%%%%%%%%
\subsection{Discrete-Time State Space}
\label{subsec:discrete-time-state-space}
%%%%%%%%%%%%%%%%%%%%%%%%%%%%%%%%%%%%%%%%%%%%%%%%%%%%%%%%%%%%

A discrete-time linear system is

\[
\boxed{
\mathbf{x}_{k+1}
=
A\mathbf{x}_k
+
B\mathbf{u}_k,
}
\]

\[
\boxed{
\mathbf{y}_k
=
C\mathbf{x}_k
+
D\mathbf{u}_k.
}
\]

Such models arise naturally from sampled digital control.

%%%%%%%%%%%%%%%%%%%%%%%%%%%%%%%%%%%%%%%%%%%%%%%%%%%%%%%%%%%%
\subsection{Discrete-Time Stability}
\label{subsec:discrete-time-stability-control}
%%%%%%%%%%%%%%%%%%%%%%%%%%%%%%%%%%%%%%%%%%%%%%%%%%%%%%%%%%%%

For

\[
\mathbf{x}_{k+1}
=
A\mathbf{x}_k,
\]

asymptotic stability requires

\[
\boxed{
|\lambda_i|<1
}
\]

for every eigenvalue of \(A\).

Thus continuous-time stability uses the left half-plane, while discrete-time
stability uses the interior of the unit disk.

%%%%%%%%%%%%%%%%%%%%%%%%%%%%%%%%%%%%%%%%%%%%%%%%%%%%%%%%%%%%
\subsection{Worked Example: Discrete Stability}
\label{subsec:worked-discrete-stability}
%%%%%%%%%%%%%%%%%%%%%%%%%%%%%%%%%%%%%%%%%%%%%%%%%%%%%%%%%%%%

\begin{workedexamplebox}{A Stable Difference System}

Let

\[
A
=
\begin{bmatrix}
0.5 & 0\\
0 & -0.8
\end{bmatrix}.
\]

The eigenvalues are

\[
0.5
\]

and

\[
-0.8.
\]

Both satisfy

\[
|\lambda|<1.
\]

Therefore,

\begin{answerbox}
\[
\boxed{
\mathbf{x}_{k+1}=A\mathbf{x}_k
\text{ is asymptotically stable}.
}
\]
\end{answerbox}

\end{workedexamplebox}

%%%%%%%%%%%%%%%%%%%%%%%%%%%%%%%%%%%%%%%%%%%%%%%%%%%%%%%%%%%%
\subsection{Sampling}
\label{subsec:control-sampling}
%%%%%%%%%%%%%%%%%%%%%%%%%%%%%%%%%%%%%%%%%%%%%%%%%%%%%%%%%%%%

Digital controllers measure and update systems at discrete times

\[
\boxed{
t_k=kT_s,
}
\]

where \(T_s\) is the sampling period.

If sampling is too slow, important system dynamics may be missed.

Thus digital control design must account for sampling frequency.

%%%%%%%%%%%%%%%%%%%%%%%%%%%%%%%%%%%%%%%%%%%%%%%%%%%%%%%%%%%%
\subsection{Continuous-to-Discrete Conversion}
\label{subsec:continuous-discrete-control-conversion}
%%%%%%%%%%%%%%%%%%%%%%%%%%%%%%%%%%%%%%%%%%%%%%%%%%%%%%%%%%%%

For

\[
\dot{\mathbf{x}}
=
A\mathbf{x}
+
B\mathbf{u},
\]

assuming a zero-order-held input over each sample interval \(T_s\),

\[
\boxed{
\mathbf{x}_{k+1}
=
A_d\mathbf{x}_k
+
B_d\mathbf{u}_k,
}
\]

where

\[
\boxed{
A_d
=
e^{AT_s},
}
\]

and

\[
\boxed{
B_d
=
\int_0^{T_s}
e^{A\tau}B\,d\tau.
}
\]

%%%%%%%%%%%%%%%%%%%%%%%%%%%%%%%%%%%%%%%%%%%%%%%%%%%%%%%%%%%%
\subsection{Time Delay}
\label{subsec:control-time-delay}
%%%%%%%%%%%%%%%%%%%%%%%%%%%%%%%%%%%%%%%%%%%%%%%%%%%%%%%%%%%%

Control systems may contain sensing, communication, computation, or actuator
delays.

A delayed input may appear as

\[
\boxed{
\dot{\mathbf{x}}(t)
=
A\mathbf{x}(t)
+
B\mathbf{u}(t-\tau).
}
\]

Delays introduce additional phase lag and can destabilize feedback systems.

%%%%%%%%%%%%%%%%%%%%%%%%%%%%%%%%%%%%%%%%%%%%%%%%%%%%%%%%%%%%
\subsection{Nonlinear Control}
\label{subsec:nonlinear-control}
%%%%%%%%%%%%%%%%%%%%%%%%%%%%%%%%%%%%%%%%%%%%%%%%%%%%%%%%%%%%

Many real systems are nonlinear:

\[
\boxed{
\dot{\mathbf{x}}
=
f(\mathbf{x},\mathbf{u}).
}
\]

Nonlinear control must account for phenomena such as:

\[
\boxed{
\text{multiple equilibria},
}
\]

\[
\boxed{
\text{saturation},
}
\]

\[
\boxed{
\text{limit cycles},
}
\]

and

\[
\boxed{
\text{state-dependent dynamics}.
}
\]

Linearization may only describe local behavior.

%%%%%%%%%%%%%%%%%%%%%%%%%%%%%%%%%%%%%%%%%%%%%%%%%%%%%%%%%%%%
\subsection{Feedback Linearization}
\label{subsec:feedback-linearization}
%%%%%%%%%%%%%%%%%%%%%%%%%%%%%%%%%%%%%%%%%%%%%%%%%%%%%%%%%%%%

Feedback linearization seeks a nonlinear control transformation that cancels
certain nonlinear terms.

For a scalar system

\[
\dot{x}
=
f(x)+g(x)u,
\]

if \(g(x)\neq0\), choose

\[
\boxed{
u
=
\frac{
v-f(x)
}{
g(x)
}.
}
\]

Then

\[
\boxed{
\dot{x}=v.
}
\]

This method requires sufficiently accurate model knowledge.

%%%%%%%%%%%%%%%%%%%%%%%%%%%%%%%%%%%%%%%%%%%%%%%%%%%%%%%%%%%%
\subsection{Adaptive Control}
\label{subsec:adaptive-control}
%%%%%%%%%%%%%%%%%%%%%%%%%%%%%%%%%%%%%%%%%%%%%%%%%%%%%%%%%%%%

Adaptive control modifies controller parameters while the system operates.

This is useful when system parameters are unknown or slowly changing.

Schematically,

\[
\boxed{
\text{controller}
+
\text{online parameter update}.
}
\]

Stability of the adaptation law is a central concern.

%%%%%%%%%%%%%%%%%%%%%%%%%%%%%%%%%%%%%%%%%%%%%%%%%%%%%%%%%%%%
\subsection{System Identification}
\label{subsec:system-identification-control}
%%%%%%%%%%%%%%%%%%%%%%%%%%%%%%%%%%%%%%%%%%%%%%%%%%%%%%%%%%%%

System identification estimates mathematical models from observed input and
output data.

A general goal is to infer

\[
\boxed{
\text{model structure}
}
\]

and

\[
\boxed{
\text{parameters}
}
\]

from experiments.

Identification and control are strongly connected because controller quality
depends on model quality.

%%%%%%%%%%%%%%%%%%%%%%%%%%%%%%%%%%%%%%%%%%%%%%%%%%%%%%%%%%%%
\subsection{Persistent Excitation}
\label{subsec:persistent-excitation}
%%%%%%%%%%%%%%%%%%%%%%%%%%%%%%%%%%%%%%%%%%%%%%%%%%%%%%%%%%%%

To estimate model parameters reliably, input signals must often excite enough
independent system behavior.

If an experiment never activates certain dynamics, those dynamics may be
unidentifiable.

This motivates the idea of sufficiently informative or persistently exciting
inputs.

%%%%%%%%%%%%%%%%%%%%%%%%%%%%%%%%%%%%%%%%%%%%%%%%%%%%%%%%%%%%
\subsection{Controllability Gramian}
\label{subsec:controllability-gramian}
%%%%%%%%%%%%%%%%%%%%%%%%%%%%%%%%%%%%%%%%%%%%%%%%%%%%%%%%%%%%

For a stable continuous-time LTI system, the infinite-horizon controllability
Gramian is

\[
\boxed{
W_c
=
\int_0^\infty
e^{At}
BB^T
e^{A^Tt}
\,dt.
}
\]

It satisfies

\[
\boxed{
AW_c
+
W_cA^T
+
BB^T
=
0.
}
\]

The Gramian quantifies how strongly different state directions can be reached
by inputs.

%%%%%%%%%%%%%%%%%%%%%%%%%%%%%%%%%%%%%%%%%%%%%%%%%%%%%%%%%%%%
\subsection{Observability Gramian}
\label{subsec:observability-gramian}
%%%%%%%%%%%%%%%%%%%%%%%%%%%%%%%%%%%%%%%%%%%%%%%%%%%%%%%%%%%%

The observability Gramian is

\[
\boxed{
W_o
=
\int_0^\infty
e^{A^Tt}
C^TC
e^{At}
\,dt.
}
\]

It satisfies

\[
\boxed{
A^TW_o
+
W_oA
+
C^TC
=
0.
}
\]

It measures how strongly different state directions appear in the output.

%%%%%%%%%%%%%%%%%%%%%%%%%%%%%%%%%%%%%%%%%%%%%%%%%%%%%%%%%%%%
\subsection{Balanced Realizations}
\label{subsec:balanced-realizations}
%%%%%%%%%%%%%%%%%%%%%%%%%%%%%%%%%%%%%%%%%%%%%%%%%%%%%%%%%%%%

A balanced realization transforms coordinates so controllability and
observability Gramians are equal and diagonal:

\[
\boxed{
W_c
=
W_o
=
\operatorname{diag}
(\sigma_1,\ldots,\sigma_n).
}
\]

The values \(\sigma_i\) are Hankel singular values.

Small values identify states that are simultaneously difficult to control
and observe.

%%%%%%%%%%%%%%%%%%%%%%%%%%%%%%%%%%%%%%%%%%%%%%%%%%%%%%%%%%%%
\subsection{Balanced Truncation}
\label{subsec:balanced-truncation}
%%%%%%%%%%%%%%%%%%%%%%%%%%%%%%%%%%%%%%%%%%%%%%%%%%%%%%%%%%%%

Balanced truncation removes state directions associated with small Hankel
singular values.

This produces a lower-order approximation while preserving important
input-output behavior.

It connects control theory with model reduction.

%%%%%%%%%%%%%%%%%%%%%%%%%%%%%%%%%%%%%%%%%%%%%%%%%%%%%%%%%%%%
\subsection{Distributed Control}
\label{subsec:distributed-control}
%%%%%%%%%%%%%%%%%%%%%%%%%%%%%%%%%%%%%%%%%%%%%%%%%%%%%%%%%%%%

Large systems may contain multiple controllers.

Each controller may use only local information.

Examples include:

\[
\boxed{
\text{vehicle formations},
}
\]

\[
\boxed{
\text{power grids},
}
\]

\[
\boxed{
\text{sensor networks},
}
\]

and

\[
\boxed{
\text{multi-agent robotics}.
}
\]

This leads to distributed and networked control.

%%%%%%%%%%%%%%%%%%%%%%%%%%%%%%%%%%%%%%%%%%%%%%%%%%%%%%%%%%%%
\subsection{Consensus Dynamics}
\label{subsec:consensus-dynamics}
%%%%%%%%%%%%%%%%%%%%%%%%%%%%%%%%%%%%%%%%%%%%%%%%%%%%%%%%%%%%

A basic continuous-time consensus model is

\[
\boxed{
\dot{x}_i
=
\sum_j
a_{ij}
(x_j-x_i).
}
\]

Using graph Laplacian \(L\),

\[
\boxed{
\dot{\mathbf{x}}
=
-L\mathbf{x}.
}
\]

Under suitable connectivity conditions, the states approach a common value.

%%%%%%%%%%%%%%%%%%%%%%%%%%%%%%%%%%%%%%%%%%%%%%%%%%%%%%%%%%%%
\subsection{Control-Theory Workflow}
\label{subsec:control-theory-workflow}
%%%%%%%%%%%%%%%%%%%%%%%%%%%%%%%%%%%%%%%%%%%%%%%%%%%%%%%%%%%%

A practical control-design workflow is:

\begin{enumerate}
    \item define the physical or mathematical system;
    \item identify control objectives;
    \item select state variables;
    \item identify available actuators;
    \item identify measurable outputs;
    \item construct the dynamical model;
    \item identify equilibrium points;
    \item linearize if appropriate;
    \item analyze open-loop stability;
    \item test controllability;
    \item test observability;
    \item identify disturbances and uncertainty;
    \item specify performance requirements;
    \item specify actuator and state constraints;
    \item choose a control architecture;
    \item design feedback gains or an optimization-based controller;
    \item design an observer when states are not measured;
    \item analyze closed-loop stability;
    \item evaluate transient response;
    \item evaluate steady-state error;
    \item evaluate disturbance rejection;
    \item evaluate robustness;
    \item test actuator saturation;
    \item account for delays and sampling;
    \item simulate nonlinear and uncertain models;
    \item verify controller implementation;
    \item validate performance experimentally;
    \item revise the model or controller when necessary.
\end{enumerate}

%%%%%%%%%%%%%%%%%%%%%%%%%%%%%%%%%%%%%%%%%%%%%%%%%%%%%%%%%%%%
\subsection{Common Errors}
\label{subsec:control-theory-common-errors}
%%%%%%%%%%%%%%%%%%%%%%%%%%%%%%%%%%%%%%%%%%%%%%%%%%%%%%%%%%%%

\subsubsection{Confusing State and Output}

The state contains the information needed to describe system evolution.

The output is what is measured or reported.

In general,

\[
\boxed{
\mathbf{y}
\neq
\mathbf{x}.
}
\]

%%%%%%%%%%%%%%%%%%%%%%%%%%%%%%%%%%%%%%%%%%%%%%%%%%%%%%%%%%%%
\subsubsection{Assuming All States Are Measurable}

Many physical systems have internal states that cannot be measured directly.

An observer or state estimator may be needed.

%%%%%%%%%%%%%%%%%%%%%%%%%%%%%%%%%%%%%%%%%%%%%%%%%%%%%%%%%%%%
\subsubsection{Confusing Open-Loop and Closed-Loop Stability}

A plant can be unstable in open loop but stabilized through feedback.

Conversely, poorly designed feedback can destabilize an otherwise stable
plant.

%%%%%%%%%%%%%%%%%%%%%%%%%%%%%%%%%%%%%%%%%%%%%%%%%%%%%%%%%%%%
\subsubsection{Assuming More Feedback Gain Is Always Better}

High gain can improve response in some respects but may also increase:

\[
\boxed{
\text{noise sensitivity},
}
\]

\[
\boxed{
\text{actuator saturation},
}
\]

\[
\boxed{
\text{oscillation},
}
\]

and

\[
\boxed{
\text{robustness problems}.
}
\]

%%%%%%%%%%%%%%%%%%%%%%%%%%%%%%%%%%%%%%%%%%%%%%%%%%%%%%%%%%%%
\subsubsection{Confusing Poles and Zeros}

Poles are associated primarily with system modes and stability.

Zeros describe input-output cancellation behavior.

They have different roles.

%%%%%%%%%%%%%%%%%%%%%%%%%%%%%%%%%%%%%%%%%%%%%%%%%%%%%%%%%%%%
\subsubsection{Using Pole Location Alone for Every Stability Question}

Pole tests apply directly to linear time-invariant systems.

Nonlinear systems require additional analysis, often using linearization or
Lyapunov methods.

%%%%%%%%%%%%%%%%%%%%%%%%%%%%%%%%%%%%%%%%%%%%%%%%%%%%%%%%%%%%
\subsubsection{Confusing Internal Stability and BIBO Stability}

These concepts are closely related for minimal LTI systems but are not
identical in every realization, especially when hidden unstable modes or
pole-zero cancellations are present.

%%%%%%%%%%%%%%%%%%%%%%%%%%%%%%%%%%%%%%%%%%%%%%%%%%%%%%%%%%%%
\subsubsection{Ignoring Controllability}

Pole placement cannot arbitrarily move modes that are uncontrollable.

Always check

\[
\operatorname{rank}\mathcal{C}.
\]

%%%%%%%%%%%%%%%%%%%%%%%%%%%%%%%%%%%%%%%%%%%%%%%%%%%%%%%%%%%%
\subsubsection{Ignoring Observability}

An observer cannot reconstruct modes that do not affect the measured output.

Always check

\[
\operatorname{rank}\mathcal{O}.
\]

%%%%%%%%%%%%%%%%%%%%%%%%%%%%%%%%%%%%%%%%%%%%%%%%%%%%%%%%%%%%
\subsubsection{Assuming Linearization Is Globally Valid}

Linearization is generally local.

A controller designed near one equilibrium may perform poorly far from that
state.

%%%%%%%%%%%%%%%%%%%%%%%%%%%%%%%%%%%%%%%%%%%%%%%%%%%%%%%%%%%%
\subsubsection{Ignoring Steady-State Error}

A stable system can still fail to track its reference accurately.

Stability and tracking accuracy are separate design objectives.

%%%%%%%%%%%%%%%%%%%%%%%%%%%%%%%%%%%%%%%%%%%%%%%%%%%%%%%%%%%%
\subsubsection{Using Derivative Control Without Noise Consideration}

Differentiation amplifies high-frequency measurement noise.

Practical derivative control is usually filtered.

%%%%%%%%%%%%%%%%%%%%%%%%%%%%%%%%%%%%%%%%%%%%%%%%%%%%%%%%%%%%
\subsubsection{Ignoring Integral Windup}

Integral action combined with actuator saturation can produce large transient
errors.

Anti-windup measures may be necessary.

%%%%%%%%%%%%%%%%%%%%%%%%%%%%%%%%%%%%%%%%%%%%%%%%%%%%%%%%%%%%
\subsubsection{Assuming PID Is Always Sufficient}

PID is powerful but may be inadequate for strongly nonlinear, highly
multivariable, constrained, delayed, or uncertain systems.

%%%%%%%%%%%%%%%%%%%%%%%%%%%%%%%%%%%%%%%%%%%%%%%%%%%%%%%%%%%%
\subsubsection{Ignoring Actuator Limits}

A theoretical input may require impossible forces, voltages, rates, or power.

Physical constraints must be included.

%%%%%%%%%%%%%%%%%%%%%%%%%%%%%%%%%%%%%%%%%%%%%%%%%%%%%%%%%%%%
\subsubsection{Ignoring Measurement Noise}

Feedback based directly on noisy measurements can cause excessive control
activity.

Filtering or state estimation may be necessary.

%%%%%%%%%%%%%%%%%%%%%%%%%%%%%%%%%%%%%%%%%%%%%%%%%%%%%%%%%%%%
\subsubsection{Ignoring Model Uncertainty}

A controller optimized for one exact model may perform poorly on the actual
system.

Robustness should be tested.

%%%%%%%%%%%%%%%%%%%%%%%%%%%%%%%%%%%%%%%%%%%%%%%%%%%%%%%%%%%%
\subsubsection{Confusing Feedforward and Feedback}

Feedforward acts using known desired motion or disturbances.

Feedback acts using measured error or state.

They serve complementary purposes.

%%%%%%%%%%%%%%%%%%%%%%%%%%%%%%%%%%%%%%%%%%%%%%%%%%%%%%%%%%%%
\subsubsection{Confusing LQR With Arbitrary Pole Placement}

LQR chooses feedback by minimizing a quadratic performance index.

Pole placement directly assigns desired closed-loop poles.

They are related but not equivalent design procedures.

%%%%%%%%%%%%%%%%%%%%%%%%%%%%%%%%%%%%%%%%%%%%%%%%%%%%%%%%%%%%
\subsubsection{Assuming Optimal Means Physically Best in Every Sense}

An optimal controller is optimal relative to the chosen:

\[
\boxed{
\text{model},
}
\]

\[
\boxed{
\text{cost function},
}
\]

and

\[
\boxed{
\text{constraints}.
}
\]

Changing these can change the optimal solution.

%%%%%%%%%%%%%%%%%%%%%%%%%%%%%%%%%%%%%%%%%%%%%%%%%%%%%%%%%%%%
\subsubsection{Ignoring Sampling in Digital Control}

A continuous-time controller cannot always be implemented digitally without
considering sampling and discretization.

%%%%%%%%%%%%%%%%%%%%%%%%%%%%%%%%%%%%%%%%%%%%%%%%%%%%%%%%%%%%
\subsubsection{Using Continuous-Time Stability Conditions in Discrete Time}

Continuous-time stability requires

\[
\operatorname{Re}(\lambda)<0.
\]

Discrete-time stability requires

\[
|\lambda|<1.
\]

%%%%%%%%%%%%%%%%%%%%%%%%%%%%%%%%%%%%%%%%%%%%%%%%%%%%%%%%%%%%
\subsubsection{Ignoring Delay}

Feedback delay adds phase lag and can reduce stability margins.

Even small delays may matter in fast systems.

%%%%%%%%%%%%%%%%%%%%%%%%%%%%%%%%%%%%%%%%%%%%%%%%%%%%%%%%%%%%
\subsubsection{Confusing Estimation and Control}

A state estimator attempts to infer the system state.

A controller chooses actions.

They interact but solve different problems.

%%%%%%%%%%%%%%%%%%%%%%%%%%%%%%%%%%%%%%%%%%%%%%%%%%%%%%%%%%%%
\subsubsection{Assuming Simulation Guarantees Real-World Stability}

Simulation explores selected models and scenarios.

Formal analysis, robustness testing, and experimental validation remain
important.

%%%%%%%%%%%%%%%%%%%%%%%%%%%%%%%%%%%%%%%%%%%%%%%%%%%%%%%%%%%%
\subsection{Applications}
\label{subsec:control-theory-applications}
%%%%%%%%%%%%%%%%%%%%%%%%%%%%%%%%%%%%%%%%%%%%%%%%%%%%%%%%%%%%

\begin{applicationbox}

\textbf{Aerospace.}
Control theory stabilizes aircraft, spacecraft, launch vehicles, and
autonomous flight systems.

\medskip

\textbf{Robotics.}
Feedback controls robot position, velocity, orientation, manipulation, and
trajectory tracking.

\medskip

\textbf{Automotive Systems.}
Cruise control, steering, braking, suspension, and powertrain systems use
feedback and estimation.

\medskip

\textbf{Process Control.}
Chemical plants use controllers to regulate temperature, pressure, flow,
composition, and production rates.

\medskip

\textbf{Electrical Systems.}
Control methods regulate motors, converters, generators, power grids, and
electronic circuits.

\medskip

\textbf{Biomechanics and Medicine.}
Control models appear in prosthetics, drug delivery, physiological
regulation, and medical devices.

\medskip

\textbf{Economics and Finance.}
Dynamic policy, portfolio adjustment, and resource allocation can be viewed
through optimal and stochastic control.

\medskip

\textbf{Energy Systems.}
Control coordinates batteries, renewable generation, buildings, microgrids,
and thermal systems.

\medskip

\textbf{Networks and Multi-Agent Systems.}
Distributed control coordinates fleets, communication systems, power
networks, and autonomous agents.

\end{applicationbox}

%%%%%%%%%%%%%%%%%%%%%%%%%%%%%%%%%%%%%%%%%%%%%%%%%%%%%%%%%%%%
\subsection{Exercises}
\label{subsec:control-theory-exercises}
%%%%%%%%%%%%%%%%%%%%%%%%%%%%%%%%%%%%%%%%%%%%%%%%%%%%%%%%%%%%

\begin{exercise}[1]
Define control theory.
\end{exercise}

\begin{exercise}[1]
Distinguish open-loop and closed-loop control.
\end{exercise}

\begin{exercise}[1]
Define a plant.
\end{exercise}

\begin{exercise}[1]
Define a control input.
\end{exercise}

\begin{exercise}[1]
Define a state variable.
\end{exercise}

\begin{exercise}[1]
Write the continuous-time LTI state-space model.
\end{exercise}

\begin{exercise}[1]
Define a transfer function.
\end{exercise}

\begin{exercise}[1]
Define a pole.
\end{exercise}

\begin{exercise}[1]
Define controllability.
\end{exercise}

\begin{exercise}[1]
Define observability.
\end{exercise}

\begin{exercise}[2]
Solve

\[
\dot{x}
=
-3x+2,
\qquad
x(0)=1.
\]
\end{exercise}

\begin{exercise}[2]
Find the equilibrium of

\[
\dot{x}
=
-ax+bu.
\]

Assume \(a\neq0\) and constant \(u\).
\end{exercise}

\begin{exercise}[2]
Linearize

\[
\dot{x}
=
x^2+u
\]

about an equilibrium.
\end{exercise}

\begin{exercise}[2]
Find the transfer function of

\[
\ddot{y}
+
3\dot{y}
+
2y
=
u.
\]
\end{exercise}

\begin{exercise}[2]
Find the poles of

\[
G(s)
=
\frac{1}
{s^2+3s+2}.
\]
\end{exercise}

\begin{exercise}[2]
Determine whether those poles define a stable continuous-time system.
\end{exercise}

\begin{exercise}[2]
Determine stability of

\[
A
=
\begin{bmatrix}
-1 & 0\\
0 & 2
\end{bmatrix}.
\]
\end{exercise}

\begin{exercise}[3]
Determine stability of a \(2\times2\) system from its characteristic
polynomial.
\end{exercise}

\begin{exercise}[2]
For

\[
\dot{x}
=
3x+u,
\]

find the proportional state-feedback gain \(K\) such that

\[
u=-Kx
\]

produces closed-loop pole \(-4\).
\end{exercise}

\begin{exercise}[2]
Construct the controllability matrix for a two-state LTI system.
\end{exercise}

\begin{exercise}[3]
Determine controllability of

\[
A
=
\begin{bmatrix}
0 & 1\\
0 & 0
\end{bmatrix},
\qquad
B
=
\begin{bmatrix}
1\\
0
\end{bmatrix}.
\]
\end{exercise}

\begin{exercise}[2]
Construct the observability matrix for a two-state system.
\end{exercise}

\begin{exercise}[3]
Determine observability of

\[
A
=
\begin{bmatrix}
0 & 1\\
0 & 0
\end{bmatrix},
\qquad
C
=
\begin{bmatrix}
0 & 1
\end{bmatrix}.
\]
\end{exercise}

\begin{exercise}[3]
Construct a Luenberger observer for a simple observable two-state system.
\end{exercise}

\begin{exercise}[2]
Write a proportional control law.
\end{exercise}

\begin{exercise}[2]
Write an integral control law.
\end{exercise}

\begin{exercise}[2]
Write a derivative control law.
\end{exercise}

\begin{exercise}[2]
Write the PID control law.
\end{exercise}

\begin{exercise}[3]
Explain how integral control can remove constant steady-state error.
\end{exercise}

\begin{exercise}[3]
Explain why derivative action amplifies measurement noise.
\end{exercise}

\begin{exercise}[2]
For

\[
G(s)
=
\frac{\omega_n^2}
{s^2+2\zeta\omega_ns+\omega_n^2},
\]

identify \(\omega_n\) and \(\zeta\).
\end{exercise}

\begin{exercise}[3]
Classify the response for several values of \(\zeta\).
\end{exercise}

\begin{exercise}[2]
Explain root locus conceptually.
\end{exercise}

\begin{exercise}[2]
Define frequency response.
\end{exercise}

\begin{exercise}[2]
Explain a Bode magnitude plot.
\end{exercise}

\begin{exercise}[3]
Explain gain margin and phase margin.
\end{exercise}

\begin{exercise}[2]
Define the sensitivity function

\[
S(s)=\frac{1}{1+L(s)}.
\]
\end{exercise}

\begin{exercise}[3]
Show that

\[
S+T=1
\]

for complementary sensitivity

\[
T=\frac{L}{1+L}.
\]
\end{exercise}

\begin{exercise}[2]
Define a Lyapunov function.
\end{exercise}

\begin{exercise}[3]
Use

\[
V(x)=\frac12x^2
\]

to analyze

\[
\dot{x}=-x-x^3.
\]
\end{exercise}

\begin{exercise}[3]
Derive

\[
\dot V
=
\mathbf{x}^T
(A^TP+PA)
\mathbf{x}
\]

for

\[
V=\mathbf{x}^TP\mathbf{x}.
\]
\end{exercise}

\begin{exercise}[3]
Solve a simple continuous Lyapunov equation.
\end{exercise}

\begin{exercise}[2]
Define an optimal control problem.
\end{exercise}

\begin{exercise}[2]
Write a quadratic LQR cost.
\end{exercise}

\begin{exercise}[2]
State the continuous algebraic Riccati equation.
\end{exercise}

\begin{exercise}[3]
Solve the scalar Riccati equation for

\[
\dot{x}=u,
\qquad
J=\int_0^\infty(x^2+u^2)\,dt.
\]
\end{exercise}

\begin{exercise}[3]
Explain how \(Q\) and \(R\) affect LQR design qualitatively.
\end{exercise}

\begin{exercise}[2]
Define the control Hamiltonian used in Pontryagin's Maximum Principle.
\end{exercise}

\begin{exercise}[3]
Write the costate equation.
\end{exercise}

\begin{exercise}[3]
Explain the Hamilton--Jacobi--Bellman equation conceptually.
\end{exercise}

\begin{exercise}[2]
Explain model predictive control.
\end{exercise}

\begin{exercise}[3]
Explain why MPC handles constraints naturally.
\end{exercise}

\begin{exercise}[2]
Define actuator saturation.
\end{exercise}

\begin{exercise}[3]
Explain integral windup.
\end{exercise}

\begin{exercise}[3]
Explain anti-windup conceptually.
\end{exercise}

\begin{exercise}[2]
Define robust control conceptually.
\end{exercise}

\begin{exercise}[3]
Give three sources of uncertainty relevant to controller design.
\end{exercise}

\begin{exercise}[2]
Write a stochastic state-space model.
\end{exercise}

\begin{exercise}[2]
Explain the role of a Kalman filter.
\end{exercise}

\begin{exercise}[3]
Explain the Kalman innovation or measurement residual.
\end{exercise}

\begin{exercise}[2]
Write the discrete-time state equation.
\end{exercise}

\begin{exercise}[2]
State the discrete-time eigenvalue stability condition.
\end{exercise}

\begin{exercise}[3]
Determine stability of a matrix with eigenvalues

\[
0.9,
\qquad
-1.1.
\]
\end{exercise}

\begin{exercise}[3]
Derive

\[
A_d=e^{AT_s}
\]

for sampled zero-input dynamics.
\end{exercise}

\begin{exercise}[3]
Explain why time delay can destabilize a feedback system.
\end{exercise}

\begin{exercise}[2]
Explain nonlinear feedback linearization conceptually.
\end{exercise}

\begin{exercise}[3]
Apply feedback linearization to

\[
\dot{x}
=
x^2+u.
\]
\end{exercise}

\begin{exercise}[2]
Define adaptive control conceptually.
\end{exercise}

\begin{exercise}[2]
Define system identification.
\end{exercise}

\begin{exercise}[3]
Explain persistent excitation conceptually.
\end{exercise}

\begin{exercise}[3]
Explain controllability and observability Gramians conceptually.
\end{exercise}

\begin{exercise}[3]
Explain balanced truncation.
\end{exercise}

\begin{exercise}[2]
Write a consensus model on a graph.
\end{exercise}

\begin{exercise}[3]
Explain the role of the graph Laplacian in consensus.
\end{exercise}

\begin{exercise}[4]
Design state feedback for a controllable two-state system using pole
placement.
\end{exercise}

\begin{exercise}[4]
Construct a full-order observer and analyze its error dynamics.
\end{exercise}

\begin{exercise}[4]
Study the separation principle for an observer-based controller.
\end{exercise}

\begin{exercise}[4]
Design a PID controller for a second-order plant.
\end{exercise}

\begin{exercise}[4]
Analyze a closed-loop system using root locus.
\end{exercise}

\begin{exercise}[4]
Construct Bode plots for a rational transfer function.
\end{exercise}

\begin{exercise}[4]
Compute gain and phase margins for a feedback loop.
\end{exercise}

\begin{exercise}[4]
Study Nyquist stability conceptually.
\end{exercise}

\begin{exercise}[4]
Construct a quadratic Lyapunov function for a stable linear system.
\end{exercise}

\begin{exercise}[4]
Use Lyapunov's direct method on a nonlinear two-state system.
\end{exercise}

\begin{exercise}[4]
Design an infinite-horizon LQR controller for a two-state system.
\end{exercise}

\begin{exercise}[4]
Derive the finite-horizon Riccati equation.
\end{exercise}

\begin{exercise}[4]
Apply Pontryagin's Maximum Principle to a minimum-energy control problem.
\end{exercise}

\begin{exercise}[4]
Derive an HJB equation for a simple deterministic control problem.
\end{exercise}

\begin{exercise}[4]
Construct a constrained model predictive controller.
\end{exercise}

\begin{exercise}[4]
Study robust stability under parameter uncertainty.
\end{exercise}

\begin{exercise}[4]
Implement a Kalman filter for a noisy state-space model.
\end{exercise}

\begin{exercise}[4]
Design an LQG controller conceptually.
\end{exercise}

\begin{exercise}[4]
Analyze a digital controller under different sampling periods.
\end{exercise}

\begin{exercise}[4]
Study stability of a delayed feedback system.
\end{exercise}

\begin{exercise}[4]
Design an anti-windup modification for a PI controller.
\end{exercise}

\begin{exercise}[4]
Study feedback linearization of a nonlinear mechanical system.
\end{exercise}

\begin{exercise}[4]
Construct an adaptive parameter estimator for a simple uncertain plant.
\end{exercise}

\begin{exercise}[4]
Perform system identification from simulated input-output data.
\end{exercise}

\begin{exercise}[4]
Compute controllability and observability Gramians numerically.
\end{exercise}

\begin{exercise}[4]
Perform balanced truncation on a stable linear model.
\end{exercise}

\begin{exercise}[4]
Analyze consensus of a networked multi-agent system.
\end{exercise}

%%%%%%%%%%%%%%%%%%%%%%%%%%%%%%%%%%%%%%%%%%%%%%%%%%%%%%%%%%%%
\subsection{Conceptual Summary}
\label{subsec:control-theory-summary}
%%%%%%%%%%%%%%%%%%%%%%%%%%%%%%%%%%%%%%%%%%%%%%%%%%%%%%%%%%%%

Control theory studies how inputs influence dynamical systems.

A standard linear state-space model is

\[
\boxed{
\dot{\mathbf{x}}
=
A\mathbf{x}
+
B\mathbf{u},
}
\]

\[
\boxed{
\mathbf{y}
=
C\mathbf{x}
+
D\mathbf{u}.
}
\]

Feedback modifies the input using measured or estimated system behavior:

\[
\boxed{
\mathbf{u}
=
-K\mathbf{x}.
}
\]

The closed-loop dynamics become

\[
\boxed{
\dot{\mathbf{x}}
=
(A-BK)\mathbf{x}.
}
\]

Controllability determines whether inputs can influence all required state
directions:

\[
\boxed{
\operatorname{rank}
\begin{bmatrix}
B & AB & \cdots & A^{n-1}B
\end{bmatrix}
=
n.
}
\]

Observability determines whether state information can be reconstructed from
outputs:

\[
\boxed{
\operatorname{rank}
\begin{bmatrix}
C\\
CA\\
\vdots\\
CA^{n-1}
\end{bmatrix}
=
n.
}
\]

Stability can be studied through eigenvalues or Lyapunov functions.

Classical control includes:

\[
\boxed{
\text{PID}
}
\]

and frequency-response methods.

Modern control includes:

\[
\boxed{
\text{state feedback},
}
\]

\[
\boxed{
\text{observers},
}
\]

\[
\boxed{
\text{LQR},
}
\]

\[
\boxed{
\text{MPC},
}
\]

\[
\boxed{
\text{robust control},
}
\]

and

\[
\boxed{
\text{stochastic control}.
}
\]

The central structure is

\[
\boxed{
\text{model}
+
\text{measurement}
+
\text{feedback}
+
\text{optimization}
}
\]

used to produce desired behavior despite uncertainty and disturbances.

%%%%%%%%%%%%%%%%%%%%%%%%%%%%%%%%%%%%%%%%%%%%%%%%%%%%%%%%%%%%
\subsection{Section Connections}
\label{subsec:control-theory-connections}
%%%%%%%%%%%%%%%%%%%%%%%%%%%%%%%%%%%%%%%%%%%%%%%%%%%%%%%%%%%%

\chapterconnection{
Control Theory combines differential equations, dynamical systems, linear
algebra, optimization, probability, and computation to determine how a
system should be influenced through feedback. It builds naturally on the
physical and continuum models of Sections~\ref{sec:mathematical-physics},
\ref{sec:continuum-mechanics}, and \ref{sec:fluid-dynamics}, where many
systems are governed by differential equations but require stabilization,
tracking, estimation, or optimization. State-space methods connect directly
with linear algebra, Lyapunov methods connect with dynamical systems,
optimal control connects with Chapter~10, and stochastic estimation connects
with Chapters~7 and~8. The next section, Network Science, studies systems
whose behavior depends fundamentally on graph structure, connectivity,
centrality, diffusion, synchronization, contagion, and interactions among
many connected components.
}

%%%%%%%%%%%%%%%%%%%%%%%%%%%%%%%%%%%%%%%%%%%%%%%%%%%%%%%%%%%%
% SECTION SOURCE TRACKING
%%%%%%%%%%%%%%%%%%%%%%%%%%%%%%%%%%%%%%%%%%%%%%%%%%%%%%%%%%%%

%------------------------------------------------------------
% 13.6 Control Theory
%------------------------------------------------------------
%
% Source basis:
%   - Standard linear and nonlinear control theory.
%   - Standard state-space, frequency-domain, optimal, stochastic,
%     robust, digital, and networked control methods.
%
% Used for:
%   - open-loop and closed-loop control;
%   - plants, controllers, actuators, and sensors;
%   - state-space models;
%   - linear time-invariant systems;
%   - matrix-exponential solutions;
%   - equilibria;
%   - linearization;
%   - transfer functions;
%   - poles and zeros;
%   - continuous-time stability;
%   - BIBO stability;
%   - feedback;
%   - state feedback;
%   - pole placement;
%   - controllability;
%   - controllability matrices;
%   - observability;
%   - observability matrices;
%   - observers;
%   - Luenberger observers;
%   - separation principle;
%   - proportional control;
%   - integral control;
%   - derivative control;
%   - PID control;
%   - steady-state error;
%   - transient response;
%   - second-order systems;
%   - damping;
%   - root locus;
%   - frequency response;
%   - Bode plots;
%   - gain and phase margins;
%   - feedforward control;
%   - disturbance rejection;
%   - sensitivity functions;
%   - Lyapunov stability;
%   - Lyapunov equations;
%   - optimal control;
%   - quadratic cost;
%   - LQR;
%   - Riccati equations;
%   - Pontryagin Maximum Principle;
%   - costates;
%   - Hamilton--Jacobi--Bellman equations;
%   - model predictive control;
%   - constraints;
%   - actuator saturation;
%   - integral windup;
%   - robust control;
%   - uncertainty;
%   - stochastic control;
%   - state estimation;
%   - Kalman filters;
%   - discrete-time systems;
%   - digital control;
%   - sampling;
%   - continuous-to-discrete conversion;
%   - delay;
%   - nonlinear control;
%   - feedback linearization;
%   - adaptive control;
%   - system identification;
%   - persistent excitation;
%   - controllability and observability Gramians;
%   - balanced realization;
%   - balanced truncation;
%   - distributed control;
%   - consensus dynamics;
%   - applications and exercises.
%
% Notes:
%   - Network Science is developed next in Section 13.7.
%   - Mathematical Finance follows in Section 13.8.
%   - Mathematical Economics follows in Section 13.9.
%
%%%%%%%%%%%%%%%%%%%%%%%%%%%%%%%%%%%%%%%%%%%%%%%%%%%%%%%%%%%%